\chapter{Strafferet: skyld, lovovertrædelse og sanktion}

\chapterlead{Strafferetten giver staten adgang til at reagere med straf på bestemte handlinger. Fordi straffen kan være meget indgribende, er strafferetten omgivet af krav om klar lovhjemmel, skyld, bevis og fair proces.}

\section{Hvorfor strafferetten er særlig}

En civil sag kan ende med betaling eller ophør af en aftale. En straffesag kan føre til bøde, fængsel, frakendelse eller andre sanktioner, som griber direkte ind i frihed og omdømme. Derfor er det ikke nok, at en handling virker skadelig. Den skal være kriminaliseret på det relevante tidspunkt, og anklagemyndigheden skal bevise de nødvendige betingelser.

Straffeloven samler de almindelige regler om strafbare handlinger og sanktioner, men særlove kan også indeholde strafbestemmelser \citep{straffeloven}.

\section{Legalitetsprincippet}

Strafferetligt legalitetsprincip kan sammenfattes sådan: Ingen straf uden lov. Princippet kræver blandt andet:
\begin{itemize}
  \item at handlingen var strafbar, da den blev begået,
  \item at straffebestemmelsen er tilstrækkeligt klar,
  \item at analogisk udvidelse til skade for den tiltalte undgås, og
  \item at en strengere strafregel ikke gives tilbagevirkende kraft.
\end{itemize}

Legalitetsprincippet beskytter forudsigelighed. En borger skal kunne forstå, hvilke handlinger der kan udløse straf, selv om man ikke kan forvente at kende hele lovsamlingen uden hjælp.

\section{Den strafbare handling}

En strafferetlig analyse kan opdeles i objektive og subjektive elementer.

\begin{description}[style=nextline,leftmargin=0pt,labelindent=0pt]
  \item[Objektive elementer] Det ydre forløb: handlingen, genstanden, resultatet, tid, sted og andre forhold, som bestemmelsen kræver.
  \item[Subjektive elementer] Den tiltaltes sindelag: forsæt, uagtsomhed eller et særligt formål.
  \item[Ulovlighed] Om handlingen er dækket af en straffebestemmelse og ikke af en straffrihedsgrund.
  \item[Skyld] Om personen kan bebrejdes handlingen efter reglerne om tilregnelighed og subjektiv skyld.
\end{description}

Det er ikke nok at bevise, at et resultat indtrådte. Man skal bevise de elementer, der er nævnt i den konkrete bestemmelse. En person kan have forårsaget en skade uden at have handlet med det krævede forsæt eller den krævede uagtsomhed.

\section{Forsæt og uagtsomhed}

Forsæt handler om personens viden og vilje i forhold til de relevante omstændigheder. Forsæt kan være direkte, hvor personen ønsker resultatet, eller foreligge, når personen indser muligheden for resultatet og accepterer den. Det er ikke nok, at en risiko blot burde være indset; det kan i stedet være uagtsomhed, hvis den konkrete bestemmelse også kriminaliserer uagtsomhed. Uagtsomhed handler om, at personen burde have handlet anderledes, selv om resultatet ikke var ønsket.

Det er en fejl at slutte direkte fra resultat til skyld: ``Han ramte personen, derfor havde han forsæt.'' Den rigtige rækkefølge er at identificere, hvad personen vidste, ønskede eller burde have indset, og derefter holde det op mod bestemmelsens krav.

\section{Forsøg og medvirken}

En person kan i visse tilfælde straffes for forsøg, selv om forbrydelsen ikke blev fuldbyrdet. Forsøg kræver normalt, at personen er begyndt at udføre handlingen med det nødvendige forsæt, og at det ikke blot var forberedelse eller overvejelse.

Medvirken betyder, at en person bidrager til en lovovertrædelse, for eksempel ved at planlægge, skaffe redskaber, opmuntre eller hjælpe. Bidraget og den subjektive forbindelse til hovedhandlingen skal vurderes. At være til stede er ikke automatisk medvirken.

\section{Straffrihedsgrunde}

Nødværge, nødret, samtykke og andre grunde kan i bestemte situationer gøre en ellers strafbelagt handling straffri eller begrænse ansvaret. Betingelserne er konkrete. Nødværge kræver eksempelvis et påbegyndt eller overhængende angreb og en reaktion, der ikke går længere end nødvendigt i situationen.

En straffrihedsgrund er ikke en moralsk fribillet. Den skal dokumenteres og passe på den konkrete handling, risiko og tidsmæssige sammenhæng.

\section{Bevis og uskyldsformodning}

I en straffesag er det anklagemyndigheden, der skal bevise skyld. Den tiltalte skal ikke bevise sin uskyld. Retten skal træffe afgørelse på grundlag af en samlet vurdering af beviserne, og rimelig tvivl skal komme den tiltalte til gode.

Bevis kan være vidneforklaringer, dokumenter, digitale spor, video, tekniske undersøgelser og tiltaltes egne forklaringer. Et digitalt spor kan være stærkt, men skal stadig vurderes med hensyn til autenticitet, kontekst og alternative forklaringer.

\section{Sanktioner}

Straffen skal fastsættes efter den relevante lov og de forhold, der gør handlingen mere eller mindre alvorlig. Domstolen kan tage hensyn til blandt andet skade, fare, forsæt, gentagelse, planlægning, tilståelse og personlige forhold inden for lovens rammer.

Strafferetten har flere begrundelser: gengældelse, forebyggelse, resocialisering, beskyttelse af samfundet og markering af normer. Begrundelserne kan pege i forskellige retninger, og de forklarer ikke altid den konkrete straf alene.

\section{Eksempel: den ødelagte rude}

En person kaster en sten mod en rude, men ved ikke, om nogen befinder sig bag den. Ruden går i stykker, og en person bliver ramt af glasskår. Analysen må adskille:
\begin{itemize}
  \item handlingen: kastet med stenen,
  \item den beskadigede ting,
  \item personskaden,
  \item hvad personen vidste eller burde indse,
  \item årsagssammenhængen, og
  \item eventuelle straffrihedsgrunde.
\end{itemize}

At én handling har flere konsekvenser, betyder ikke nødvendigvis, at alle konsekvenser er omfattet af samme forsæt eller samme bestemmelse.

\caution{Strafferetlige eksempler i en lærebog er ikke en måde at diagnosticere, om en bestemt person har begået en forbrydelse. Den konkrete vurdering afhænger af den gældende bestemmelse, beviset og procesmaterialet.}

\question{En person downloader et program, der kan bruges både lovligt og ulovligt, og programmet findes senere på en computer, hvor der er begået et indbrud i et system. Hvilke elementer kan programmet bevise - og hvilke kan det ikke bevise alene?}

\section*{Kapitelopsamling}
\begin{itemize}
  \item Strafferetten kræver klar lovhjemmel, en strafbar handling og den krævede skyld.
  \item Resultatet af en handling beviser ikke automatisk forsæt eller uagtsomhed.
  \item Forsøg og medvirken kræver særskilte vurderinger.
  \item Straffrihedsgrunde og bevisregler kan være afgørende for resultatet.
  \item Uskyldsformodningen og bevisbyrden beskytter den tiltalte mod straf på utilstrækkeligt grundlag.
\end{itemize}
